\section{Gaussian restrictions of the direction moments}
\label{sec:gaussian-moments}
We strengthen the counting argument by retaining zeros at which a short
prime sum is positive. This supplies a lower bound for the derivative
on a set that still has both strict velocity signs. Throughout this
section the fixed mollifier, its constants $a,b$, and the phase
$\psi_\chi$ are those of Section~\ref{sec:signed-density}. Write
$N=N_{\chi,T}$, $L=\log T$, $V=\log L$, and put
\[
 d_T=e^{-(\log V)^2},\quad z_T=d_T^{-1},\quad x_T=T^{d_T},\quad
 S_T=\sum_{\substack{z_T<p\le x_T\\p\nmid q}}\frac1p,
 \quad \sigma_T^2=S_T/2.
\]
The prime number theorem, followed by partial summation, gives
\[
 S_T=V-(\log V)^2-2\log\log V+O_\chi(1)\sim V.
\]
Only this asymptotic is needed. It also follows from the uniform
character-growth estimate \eqref{eq:uniform-character-growth} by Perron
applied to $-\zeta'/\zeta$: shifting to $1/2+\epsilon$ gives
$\sum_{n\le u}\Lambda(n)=u+O_\epsilon(u^{1/2+2\epsilon})$;
partial summation and the convergent higher-prime-power contribution
then give $\sum_{p\le u}p^{-1}=\log\log u+C+o(1)$.
Define
\begin{align*}
 P_T(s)&=\sum_{z_T<p\le x_T}\chi(p)p^{-s},&
 X_\rho&=\Re P_T(\rho)/\sigma_T,\\
 A_\rho&=L_\chi'(\rho)B_{\chi,T}(\rho)^2,&
 C_\rho&=K_\rho|B_{\chi,T}(\rho)|^2.
\end{align*}
Here $K_\rho$ is the real quantity in \eqref{eq:square-cancel}.
Let $m_k$ denote the standard Gaussian moments.

\begin{theorem}\label{thm:gaussian-direction-moments}
For each fixed nonnegative integer $k$, as $T\to\infty$,
\begin{align*}
 N^{-1}\sum X_\rho^k&\longrightarrow m_k,&
 (NL)^{-1}\sum A_\rho X_\rho^k&\longrightarrow a m_k,\\
 (NL^2)^{-1}\sum|A_\rho|^2X_\rho^k&\longrightarrow b m_k,&
 N^{-1}\sum e^{2i\psi_\chi(\gamma)}X_\rho^k&\longrightarrow0,\\
 (NL)^{-1}\sum C_\rho X_\rho^k&\longrightarrow0.
\end{align*}
Every sum counts all zero occurrences in $(T,2T]$. The same five
limits hold with $X_\rho^k$ replaced by $\mathbf1_I(X_\rho)$,
and $m_k$ replaced by the Gaussian mass of any fixed finite interval $I$.
\end{theorem}

\subsection{Uniform rational sampling}
We record the two sampling calculations needed in the proof. They
retain the multiplicities and the arithmetic coefficients.
At endpoints a bounded distance from $T,2T$ and a distance
$\gg_\chi L^{-1}$ from the ordinates, the contour for $L_\chi'/L_\chi$
gives, for rational $x>1$ of polynomial size,
\begin{equation}\label{eq:rational-landau-new}
 \sum_{T<\gamma\le2T}x^\rho
 =-\frac T{2\pi}\chi(x)\Lambda(x)
 +O_\chi\left(xL^2+
       \frac{x\log(2x)}{\operatorname{dist}_*(x,\mathbb Z)}
       +\frac L{\log x}\right).
\end{equation}
The main coefficient is zero unless $x$ is an integer prime power;
at an integer the distance means the distance to the other integers.
To prove the estimate, integrate $(L_\chi'/L_\chi)(s)x^s$ with right
edge $1+1/L$. The term $n=x$ of its absolutely convergent series is
the displayed main term. The closest remaining integer costs
$O(x\log(2x)/\operatorname{dist}_*)$; splitting the other integers
into $n<x/2$, $x/2\le n\le2x$ and $n>2x$, and using
$|\log(n/x)|\gg |n-x|/x$ in the middle range, gives $O(xL^2)$.
On the reflected left edge, integration by parts in the gamma
logarithm gives $O(L/\log x)$; the absolutely convergent reflected
series is bounded by the same estimate. The good horizontal edges
give $O(xL^2)$ by partial fractions and the local zero count.
These estimates prove \eqref{eq:rational-landau-new}. Moving the
endpoints changes $O_\chi(L)$ summands, and will always be estimated
separately below.

We also need the even Hardy phase at rational frequencies. For coprime
$a,b\le T$ and $T^{-1}\le b/a\le T$, its uniform form is
\begin{align}\label{eq:rational-phase-new}
 \sum_{0<\gamma\le T}e^{2i\theta_\chi(\gamma)}(b/a)^{i\gamma}
 ={}&-\frac{qT}{2\pi\varphi(q)}\sqrt{b/a}\,
       \frac{\chi(b)\overline{\chi(a)}\mu(b)}{\varphi(b)}\nonumber\\
 &+O_{\chi,\epsilon}\left(T^{1/2+\epsilon}\sqrt b+
                         (a/b+\sqrt{a/b})L^2\right).
\end{align}
Here any fixed $\epsilon>0$ is allowed. A direct proof keeps the
weights in \eqref{eq:twisted-hardy}. Move its left edge to $-1/2$.
Stirling and the functional equation bound that edge by
$O_\chi(x^{-1}L^2)$ and the top by $O_\chi(\sqrt{xT}L^2)$.
In the stationary-phase calculation isolate the nearest integers to
$qxT/(2\pi)$; the bound
\[
 \sum_{n\asymp B}\frac1{|n-B|+b_0\sqrt T}
 \ll\frac1{b_0\sqrt T}+1+\log(2+T),\qquad b_0=qx,
\]
gives the remaining $x^{-1/2}L^2$ error. Thus the prime sum in
\eqref{eq:twisted-hardy} is uniform in this whole range, with error
$O_\chi((\sqrt{xT}+x^{-1}+x^{-1/2})L^2)$.

For $x=b/a$, let $Q=qb$, $Y=qbT/(2\pi a)$. The exact induced Gauss
sum is
\[
 \sum_{\substack{u\bmod Q\\(u,Q)=1}}\chi(u)e^{-2\pi iau/Q}
       =\mu(b)\chi(b)\overline{\chi(-a)}\tau(\chi).
\]
Indeed insert $\sum_{d\mid(u,b)}\mu(d)$, put $u=dv$ and sum over
the lifts modulo $q$. Their additive sum vanishes unless $b/d$ divides
$a$, so coprimality forces $d=b$. The remaining primitive Gauss sum
is the asserted one, including its zero value when $(ab,q)>1$.
Expand $\chi(u)e^{-2\pi iau/Q}$ in characters modulo $Q$.
Parseval says the squared sum of its coefficients is one, so their
absolute sum is at most $\sqrt{\varphi(Q)}$. For each character the
Perron argument just used, now with
\eqref{eq:uniform-character-growth} and Cauchy's formula for its
logarithmic derivative, gives
\[
 \sum_{n\le Y}\psi(n)\Lambda(n)
   =\mathbf1_{\psi\ {\mathrm{principal}}}Y
          +O_\epsilon(Y^{1/2+\epsilon}Q^\epsilon).
\]
Extra prime powers for an induced character cost only
$O(\omega(Q)\log(2Y))$. Apply Parseval, divide by $\sqrt{qb/a}$,
and decrease the chosen growth exponent if necessary. The resulting
error is the first term in \eqref{eq:rational-phase-new}.
The root number $\tau(\chi)/(i^{e_\chi}\sqrt q)$ cancels the Gauss
phase. The remaining main coefficient is the one displayed there.
When $(b,q)>1$ it vanishes; otherwise
$\varphi(qb)=\varphi(q)\varphi(b)$. This proves all cases.

\subsection{The first moment with two coefficient polynomials}
For the prime insertions we need a bilinear version of the first
moment, which we state explicitly. Let
$X_\chi(s)=\sum x(n)\chi(n)n^{-s}$ and
$Y_{\bar\chi}(s)=\sum y(n)\bar\chi(n)n^{-s}$ have real divisor-bounded
coefficients and length $y\le T^\theta$, with fixed $\theta<1/2$.
The moment $\sum_{0<\gamma\le U}L_\chi(\rho+\alpha)
X_\chi(\rho)Y_{\bar\chi}(1-\rho)$ equals, up to
$O_{q,\epsilon}(T^{1/2+\epsilon}y)$,
\begin{align}\label{eq:bilinear-first-shift}
 &\frac1{2\pi}\int_0^U\log\frac{qt}{2\pi}
       \sum_{(n,q)=1}\frac{(n^{-\alpha}*x)(n)y(n)}n\,dt
 -\frac U{2\pi}\sum_{(n,q)=1}
       \frac{(\Lambda*n^{-\alpha}*x)(n)y(n)}n\nonumber\\
 &\quad+\frac1{2\pi}\int_0^U\left.\partial_\gamma
       \sum_{\substack{(h,k)=1\\(ghk,q)=1}}
       \frac{y(gh)x(gk)}{ghk}
       \left\{Z^q_{\alpha,\gamma,h,k}
       +(qt/(2\pi))^{-\alpha-\gamma}
                    Z^q_{-\gamma,-\alpha,h,k}\right\}
                              \right|_{\gamma=0}dt,
\end{align}
where
\[
 Z^q_{\alpha,\gamma,h,k}
 =h^{-\alpha}k^{-\gamma}
   \frac{Z_q(1+\alpha+\gamma)}{Z_q(1+\alpha)}
        \prod_{p\mid k}\frac{1-p^\gamma}{1-p^{-1-\alpha}}.
\]
Shifts are $O(L^{-1})$ and coinciding-pole values are removable.
For zeta these are exactly the two equivalent published main terms
of Benli--Elma--Ng \cite[Theorem~1.1, equations~(9)--(14),
pp.~3--4]{benli-elma-ng}. The divisor-bounded version used here,
including the coefficients with prime insertions, follows from the
contour derivation below and Lemma~\ref{lem:dirichlet-additive}.

Here is the contour derivation for the fixed character, using the
already proved error machinery. Integrate
$(L_\chi'/L_\chi)(s)L_\chi(s+\alpha)X_\chi(s)Y_{\bar\chi}(1-s)$.
The right-edge diagonal is the negative convolution term above, since
its indices satisfy $nh=k$ and their character factors cancel.
The harmonic off-diagonal estimate in \eqref{eq:real-bilinear} bounds
the remainder. Reflect the left edge. Before the outside minus sign,
its upper-right expression is the conjugate of
\[
 \frac1{2\pi i}\int
 \Delta_\chi(s-\bar\alpha)^{-1}
 \left(\frac{\Delta_\chi'}{\Delta_\chi}(s)
                         -\frac{L_\chi'}{L_\chi}(s)\right)
 L_\chi(s-\bar\alpha)Y_\chi(s)X_{\bar\chi}(1-s)\,ds.
\]
Replace the parentheses temporarily by $L_\chi(s+\gamma)/L_\chi(s)$,
rename the barred shift, and put $v=s-\alpha$. The inner convolution
is $1*n^{-\alpha-\gamma}*(\mu(n)n^{-\alpha})$.
Lemma~\ref{lem:dirichlet-additive} applies with three factors and the
two different coefficient polynomials. Its principal series is
$Z_q(s)Z_q(s+\alpha+\gamma)/Z_q(s+\alpha)$.
At $p^v\parallel k$, the same principal weights as in that lemma are
zero below $v-1$, $-1/(p-1)$ at $v-1$, and one thereafter.
The local numerator has coefficients $f_0=1$ and
$f_j=(1-p^\gamma)p^{-j(\alpha+\gamma)}$ for $j\ge1$.
Summing its tail gives the local ratio
\[
 \frac{(1-p^\gamma)p^{-v(1+\alpha+\gamma)}}{1-p^{-1-\alpha}}.
\]
Restoring the outer cutoff and the $h$ factor gives $Z^q$.
The residue $\prod_{p\mid q}(1-p^{-1})$ cancels $q/\varphi(q)$.
The pole $1-\alpha-\gamma$ gives the other term by the substitution
$(\alpha,\gamma)\mapsto(-\gamma,-\alpha)$.
At $\gamma=0$ the first term has only $k=1$, while the second
vanishes through $1/Z_q(1)$. This gives the diagonal with
$n^{-\alpha}*x$ when the gamma logarithm is restored.
Differentiation in $\gamma$ gives the last term of
\eqref{eq:bilinear-first-shift}, with its positive sign.
The horizontal and Perron errors are those already proved in
Section~\ref{sec:dirichlet-moments}; uniform Cauchy differentiation
adds only fixed logarithmic powers. This proves the interface.

\subsection{Proof of the five limits}
\begin{proof}[Proof of Theorem~\ref{thm:gaussian-direction-moments}]
First expand $P_T(\rho)^j\overline{P_T(\rho)}^{\,l}$ for fixed
$j,l$. Equal products of primes require $j=l$, and unique
factorization gives the diagonal $j!S_T^j+O_j((1+S_T)^{j-1}/z_T)$.
For the off-diagonal main terms in \eqref{eq:rational-landau-new},
an integer prime-power ratio adds copies of just one prime.
Their absolute contribution is at most
$C_{j,l}T(1+V)^{j+l}\sum_{p\le x_T}\log p/p
\ll T d_TL(1+V)^{j+l}$.
After normalization this tends to zero, since $d_TV^r\to0$ for
every fixed $r$. All remaining indices are $T^{o(1)}$ and their
nonzero rational distances and logarithms are $T^{-o(1)}$.
The contour errors, even summed absolutely, are
$T^{o(1)}L^{O_{j,l}(1)}=o(N)$, as are the endpoint strips.
Expanding the real part proves the unweighted Gaussian moments.

For the two mollifier moments, fix $K\ge k$, work through degree $K$
in a formal variable $z$, and put
\[
 E_K(w)=\sum_{j=0}^K w^j/j!,\qquad
 Q_z(s)=B_{\chi,T}(s)^2E_K(zP_T(s)/(2\sigma_T)).
\]
Use $X=Q_z$, $Y=E_K(zP_T/(2\sigma_T))$ in
\eqref{eq:bilinear-first-shift}, and use $Q_z$ on both sides in
\eqref{eq:dirichlet-discrete}. The coefficient of $z^k$ is the
desired insertion divided by $k!$. The largest length is
$T^{2\nu+Kd_T}<T^\theta$ for one fixed $\theta<1/2$, eventually.
The coefficients are uniformly divisor-bounded: a $j$th prime
power has coefficient at most $j!\tau_j(n)$, and
$\sigma_T^{-1}\le1$. Interpolation at $2K+1$ real bounded values of
$z$ supplies each coefficient while preserving all error bounds.

We give the Euler-product and tail checks. At a marked prime introduce
insertion variables $u,v$, and put $x=-1+u$, $y=-1+v$ for the
two squared-mollifier coefficients. The coefficients of $p^{-1}$ in
the four relevant products, with all scaled shifts zero, are
\[
 \begin{array}{c|c}
 \text{first diagonal}&v(1+x)=uv\\
 \text{first two-residue product}&xv+v=uv\\
 \text{second gcd product}&xy+x+y=-1+uv\\
 \text{second full three-residue product}&1+xy+x+y=uv.
 \end{array}
\]
These follow by listing $g=p$, $h=p$, $k=p$ in the exact gcd sums.
In the first two-residue product the $k$-only term vanishes through
$1-p^\gamma$; the final row includes its explicit zeta prefactor.
Thus every ratio to its uninserted product has first coefficient
$uv/p$. Higher coefficients through fixed formal degree are
$O_K(p^{-2})$. The uninserted first two-residue product itself is
\[
 C_1\frac{Z_q(1+\alpha+\gamma)}{Z_q(1+\alpha)}
       \prod_{i=1}^2
       \frac{Z_q(1+u_i)^{1/2}}{Z_q(1+u_i+\gamma)^{1/2}},
\]
where $C_1$ converges absolutely near zero and equals one there.
The two second-moment products are
\eqref{eq:square-euler-F}--\eqref{eq:square-euler-G}.

Use the scaled Mellin weights from the square-moment proof. Their
decay is $O((1+|\Im s|)^{-J-1})$, $J\ge16$.
On $|\Im s_i|\le R=d_T^{-1/2}$, replacing the local shifted powers
by one costs $O(d_TR)$ per marked first coefficient. Their product,
with $u=v=z/(2\sigma_T)$, is therefore, coefficientwise,
\[
 \frac{\text{inserted product}}{\text{uninserted product}}
    =\exp(z^2S_T/(4\sigma_T^2))+o(1)=e^{z^2/2}+o(1).
\]
The error is a fixed power of $1+V$ times
$d_TRV+z_T^{-1}$, which tends to zero. This is uniform on separated
physical shift circles of radius $O(L^{-1})$.
On all Mellin lines the insertion coefficients are $O_K((1+V)^K)$.
The original global bounds for the second-moment products are
$O_q(L^3)$ and $O_q(L^6)$, and the first two-residue product is
$O_q(L^2)$. Removing $|\Im s_i|>\sqrt L$ costs respectively
$O(V^KL^{3-J/2})$, $O(V^KL^{6-J/2})$ and
$O(V^KL^{2-J/2})$, negligible at the required scales.
On the intermediate region $R<\max|\Im s_i|\le\sqrt L$ the Laurent
approximation gives normalized kernels bounded by
$C\prod_i(1+|\Im s_i|)^2$. Its integral is
$O(V^KR^{-14})=o(1)$. The absolutely convergent Euler remainders
tend uniformly to one in the central box.

The first diagonal has the same Gaussian factor and uninserted value
one. In its $\Lambda$ convolution every index supplied by $Y$ is a
product of at most $K$ marked primes, so its coefficient is
$O_K(\log x_T(1+V)^K)=o(L)$; it remains negligible after physical
shift differentiation. The errors $T^{1/2+\epsilon}T^{2\nu+Kd_T}$
are $o(T)$ for a fixed sufficiently small $\epsilon>0$.
All conductor residues cancel as in the square-moment proof.
Cauchy differentiation in the physical shifts and extraction of
$z^k/k!$ prove the limits with weights $A_\rho$ and $|A_\rho|^2$.

For the even phase expand the exterior Euler ratio and truncate its
two indices at $M=T^{1/2}$. Its tail is $O(M^{-1/2})$, while the
fixed prime insertion is $T^{o(1)}$, so this costs $o(N)$.
The arctangent correction also costs $o(N)$.
Apply \eqref{eq:rational-phase-new} to each resulting reduced ratio.
All its numerator and denominator indices are at most $MT^{o(1)}$.
Here is an absolute bound for the main coefficient. At a prime let
$i,j$ be its exponents in the two inserted products and $r,s$ its
exponents in the exterior numerator and denominator. After the
square-root frequency factor and $1/\varphi(b)$ are included, its
absolute weight is at most
\[
                    2p^{-\max(i+r,j+s)-r-s}.
\]
At an unmarked prime the sum is $1+O(p^{-2})$; at a marked prime
it is $O_k(p^{-\max(i,j)})$. Hence the full main coefficient is
$O_k((1+V)^k)$, giving $O_{\chi,k}(T(1+V)^k)=o(N\sigma_T^k)$.
The $\sqrt b$ error has only one harmonic exterior sum, since
$\sqrt b\le\sqrt{m h}$ with $h$ an inserted prime product.
The other exterior sums converge or cost $O(\sqrt M)$.
Thus, for fixed sufficiently small $\epsilon$, all errors are
$T^{1/2+\epsilon+o(1)}=o(N)$. This proves the phase limit.

For the real signed weight expand $X_\rho^k$ into its monomials.
Apply \eqref{eq:real-bilinear} to $BP_T^j$ and $BP_T^{k-j}$, with
the exterior series truncated at $T^{\delta_0}$, where
$\nu+\delta_0<1/2$. The imaginary part is $o(NL)$.
Restoring the exterior tail costs at most
\[
 C T^{-\delta_0/2}
 \left(\sum|A_\rho|^2\sum|X_\rho|^{2k}\right)^{1/2}=o(NL).
\]
The additional real term divided by $2\gamma$ in the exact formula
for $K_\rho$ is $O(NL/T)$ by the same inequality. This proves the
fifth polynomial-moment limit.

Finally consider the positive measures of masses $1/N$ and
$|A_\rho|^2/(NL^2)$ at $X_\rho$. Taylor's formula for their
characteristic functions has limiting remainder at degree $2j$
bounded by $C|t|^{2j}m_{2j}/(2j)!=C|t|^{2j}/(2^jj!)$.
Their characteristic functions therefore tend to those of Gaussian
measures of masses one and $b$. Fourier inversion gives convergence
on smooth compactly supported functions; upper and lower smooth
approximations give convergence on intervals.
For the complex first measure, the Taylor remainder is bounded by
\[
 \frac{|t|^k}{k!}
 \left(\frac1{NL^2}\sum|A_\rho|^2\right)^{1/2}
 \left(\frac1N\sum|X_\rho|^{2k}\right)^{1/2},
\]
whose limiting bound tends to zero. Its variation on any set is
bounded by the geometric mean of the two positive masses there.
The phase measure is dominated by the first positive measure, and
$|C_\rho|\le6|A_\rho|$ dominates the signed measure. The same Taylor,
Fourier and boundary-approximation argument gives all three remaining
interval limits. Gaussian boundary masses vanish, completing the proof.
\end{proof}
